\documentclass[11pt]{amsart}
\usepackage{amsmath,amssymb,lmodern}
\usepackage[hidelinks]{hyperref}
\hypersetup{pdftitle={The Complexity of Minimum-Energy Strong Orientations},pdfauthor={Yitzchak Shmalo}}
\newtheorem{thm}{Theorem}[section]
\newtheorem{lem}[thm]{Lemma}
\newtheorem{cor}[thm]{Corollary}
\newtheorem{prop}[thm]{Proposition}
\theoremstyle{definition}
\newtheorem{ex}[thm]{Example}
\theoremstyle{remark}
\newtheorem{question}[thm]{Question}
\numberwithin{equation}{section}
\hyphenation{Vo-ro-noi}
\DeclareMathOperator{\dist}{dist}
\newcommand{\Emin}{E_{\min}}
\newcommand{\Z}{\mathbb Z}
\newcommand{\R}{\mathbb R}

\title[Minimum-energy strong orientations]{The Complexity of Minimum-Energy Strong Orientations}
\author[Yitzchak Shmalo]{Yitzchak Shmalo}
\address{Einstein Institute of Mathematics, The Hebrew University of Jerusalem, Jerusalem, Israel}
\email{yitzchak.shmalo@gmail.com}
\subjclass[2020]{Primary 05C20; Secondary 68Q17, 11H31, 05C85}
\keywords{strongly connected orientation, Laplacian energy, lattice of integer flows, covering radius, NP-completeness, treewidth}

\newcommand{\doi}[1]{doi:\href{https://doi.org/#1}{#1}}
\begin{document}

\begin{abstract}
For an orientation of a finite graph, let $f$ solve $\Delta f=d^{\mathrm{out}}-d^{\mathrm{in}}$ and call $\langle f,\Delta f\rangle$ the energy of the orientation. Amini showed that the covering radius of the lattice of integer flows is determined by the least energy of a strongly connected orientation, and wrote that the complexity of finding such an orientation is unknown and that he suspected it to be NP-hard. We show that the decision version is NP-complete, even for simple bipartite planar graphs that become a tree after one vertex is deleted. Hence deciding whether the squared covering radius of the lattice of integer flows is at most a given rational number is coNP-complete, and the same holds for lattices of integer cuts. On graphs of bounded treewidth the least energy has a fully polynomial-time approximation scheme.
\end{abstract}

\maketitle

\section{Introduction}

Let $G=(V,E)$ be a finite connected graph without loops; parallel edges are allowed. The lattice of integer flows $\Lambda(G)$ is the set of integer-valued circulations in $\Z^E$, with the inner product inherited from $\R^E$. Bacher, de la Harpe and Nagnibeda \cite{BHN97} studied this lattice and the lattice of integer cuts. Amini \cite[Theorem~1]{Amini} showed that the face poset of the Voronoi cell of $\Lambda(G)$ is isomorphic to the poset of strongly connected orientations of subgraphs of $G$, as conjectured by Caporaso and Viviani \cite{CV10}; Dancso and Lim \cite{DL22} extended this description to regular matroids. In particular the vertices of the Voronoi cell are points $v^D$ attached to the strongly connected orientations $D$ of $G$ \cite[Lemma~8]{Amini}, and the squared covering radius of $\Lambda(G)$ is the largest of their squared norms \cite[Proposition~12]{Amini}. Amini then expressed it through the following quantity.

For an orientation $D$ of $G$ let $d^{\mathrm{out}}_y$ and $d^{\mathrm{in}}_y$ be the out-degree and the in-degree of $y\in V$, and let $f_D\colon V\to\R$ solve the Laplacian equation $\Delta f_D=d^{\mathrm{out}}-d^{\mathrm{in}}$, where $\Delta f(y)=\sum_{z\sim y}(f(y)-f(z))$ with multiplicity. The \emph{energy} of $D$ is
\[
E(D)=\langle f_D,\Delta f_D\rangle=\sum_{y\in V}f_D(y)\bigl(d^{\mathrm{out}}_y-d^{\mathrm{in}}_y\bigr),
\]
which does not depend on the choice of the solution $f_D$. By Robbins' theorem \cite{Robbins39}, $G$ has a strongly connected orientation, also called a strong orientation, exactly when it is bridgeless; write $\Emin(G)$ for the least energy of a strongly connected orientation of a bridgeless graph $G$ with $m$ edges. Amini calls such an orientation well-balanced and expressed the squared covering radius of $\Lambda(G)$ as an affine function of $\Emin(G)$ \cite[Definition--Theorem~14]{Amini}; it equals $(m-\Emin(G))/4$ (Corollary~\ref{cor:radius}). Amini wrote that the complexity of finding a strongly connected orientation of least energy is not known, and that he suspected it to be NP-hard \cite[\S4]{Amini}.

A lattice $L\subset\R^n$ is regular if the support of every nonzero vector contains the support of an elementary vector, a nonzero vector of $L$ with entries in $\{0,\pm1\}$ and minimal support. Dutour Sikiri\'c, Sch\"urmann and Vallentin \cite[Theorem~2]{DSV09} showed that the Voronoi cell of a regular lattice is the orthogonal projection of the cube $[-1/2,1/2]^n$ onto the span of $L$. They gave a formula for the covering radius of the lattice of integer cuts of a graph \cite[Proposition~3]{DSV09} and remarked after it that a combinatorial interpretation of the formula could lead to a proof that the covering radius problem is NP-hard. Vallentin has since shown that the covering radius problem for general Euclidean lattices is complete for the second level of the polynomial hierarchy \cite{Vallentin26}.

Deciding whether $\Emin(G)\le r$ is therefore deciding whether the squared covering radius of $\Lambda(G)$ is at least $(m-r)/4$. Vallentin's theorem~\cite{Vallentin26} does not reach these lattices. $\Lambda(G)$ is the integer kernel of the totally unimodular incidence matrix, so its closest vector problem is solvable in polynomial time~\cite{MPSV21}, and that algorithm does not decide the covering radius.

Write an orientation as a sign vector $s\in\{\pm1\}^E$ and let $C$ be the orthogonal projection of $\R^E$ onto the space of circulations. Since $\Lambda(G)$ is regular, its Voronoi cell is $C([-1/2,1/2]^E)$~\cite[Theorem~2]{DSV09}, and with Amini's description of its vertices
\[
\Emin(G)=m-\max_{s\in\{\pm1\}^E}\|Cs\|^2 .
\]
Every maximizer is a strongly connected orientation (Lemma~\ref{lem:strong}), so the maximum may be taken over all orientations. On the following planar graphs it has a closed form. Given integers $0<A_i<B_i/2$ for $i=1,\dots,k$, let $H$ have two vertices $o_{\mathrm{in}}$, $o_{\mathrm{out}}$ and vertices $h_1,\dots,h_k$. Each $h_i$ is joined to $o_{\mathrm{in}}$ by a path with $2A_i$ edges and to $o_{\mathrm{out}}$ by two paths with $2(B_i-2A_i)$ edges each. Then (Proposition~\ref{prop:ring})
\[
\Emin(H)=W-4U+\frac{4\rho^2}{S},\qquad \rho=\min_{\sigma\in\{\pm1\}^k}\dist\Bigl(\sum_{i=1}^k\sigma_i\frac{A_i}{B_i},\ \Z\Bigr),
\]
where $W$ is the number of edges of $H$, $U=\sum_i(B_i-2A_i+A_i^2/B_i)$ and $S=\sum_i1/B_i$. The least energy of $H$ thus measures how close a signed sum of the fractions $A_i/B_i$ comes to an integer. Take a monotone not-all-equal formula in which every clause has three variables and every variable occurs in four clauses; deciding its satisfiability is NP-complete~\cite[Theorem~1]{DD20}. Add a new variable to each clause and give the $j$-th clause the $j$-th prime $p_j$ larger than $10$. Let $A_i/B_i$ be the sum of $1/p_j$ over the clauses containing the $i$-th variable. Then $\rho=0$ when the formula is satisfiable and $\rho\ge1/Q$ otherwise, where $Q=p_1\cdots p_q$ (Lemma~\ref{lem:signed}). Since every variable occurs in at most four clauses, $B_i\le p_q^4$, and $H$ has polynomial size.

Other objectives over orientations have been studied: minimizing the lexicographic order of the indegrees over acyclic orientations is NP-hard \cite{BIMOWZ17}, and so is minimizing $\sum_v\varphi(d^{\mathrm{in}}(v))$ over acyclic orientations for every fixed convex function $\varphi$ that is not linear on the positive integers \cite{BM25}. Both objectives are separable functions of the indegrees minimized over acyclic orientations, while $E(D)$ couples all edges through the Laplacian and is minimized over strongly connected orientations.

The decision problem is the following.
\begin{quote}
Given a bridgeless graph $G$ and a rational number $r$, is there a strongly connected orientation of $G$ of energy at most $r$?
\end{quote}

\begin{thm}\label{thm:hard}
The decision problem is NP-complete, even when $G$ is a simple bridgeless bipartite planar graph from which deleting one vertex leaves a tree. In particular it is NP-complete on graphs of treewidth two.
\end{thm}

For lattices coming from graphs, deciding whether the squared covering radius is at most $r$ lies in coNP, since a vertex of the Voronoi cell far from the origin is a short certificate, and Theorem~\ref{thm:hard} shows that it is coNP-hard for them.

\begin{cor}\label{cor:cov}
Deciding, for a graph $G$ and a rational number $r$, whether the squared covering radius of $\Lambda(G)$ is at most $r$ is coNP-complete, even for the graphs of Theorem~\ref{thm:hard}. The same holds for lattices of integer cuts of graphs, even of planar graphs.
\end{cor}

\begin{thm}\label{thm:fptas}
For each fixed $\omega$ there is an algorithm that, given a bridgeless loopless graph $G$ of treewidth at most $\omega$ with $m$ edges and a rational $0<\varepsilon\le1$, returns in time polynomial in $m$ and $1/\varepsilon$ a strongly connected orientation $D$ of $G$ with $E(D)\le(1+\varepsilon)\Emin(G)$ and a rational $\ell$ with $(1-\varepsilon)\Emin(G)\le\ell\le\Emin(G)$.
\end{thm}

The reduction of Theorem~\ref{thm:hard} separates its two cases by an amount exponentially small compared with the energies involved. Unless $\mathrm{P}=\mathrm{NP}$ this cannot be avoided at bounded treewidth: if the two cases always differed by a fraction of the energy at least the inverse of a polynomial, Theorem~\ref{thm:fptas} with $\varepsilon$ that small would decide them in polynomial time.

Section~\ref{sec:prelim} recasts the energy as a squared distance and shows that the strong connectivity constraint can be dropped. Section~\ref{sec:hard} proves Theorem~\ref{thm:hard}, by a reduction from monotone not-all-equal 3-satisfiability with four occurrences per variable \cite{DD20}, and Corollary~\ref{cor:cov}, and Section~\ref{sec:fptas} proves Theorem~\ref{thm:fptas}. Section~\ref{sec:open} lists questions that remain open.

\section{Energy as a distance}\label{sec:prelim}

Fix a reference orientation of $G$ and let $N\in\R^{E\times V}$ be the incidence matrix, with $N_{e,u}=1$ and $N_{e,v}=-1$ for $e=(u,v)$. An orientation is a sign vector $s\in\{\pm1\}^E$, with $s_e=1$ when $e$ keeps its reference direction. Then $N^{\top}s=d^{\mathrm{out}}-d^{\mathrm{in}}$ and $N^{\top}N=\Delta$. Let $\Delta^{+}$ be the Moore--Penrose pseudoinverse of $\Delta$, let $P=N\Delta^{+}N^{\top}$ be the orthogonal projection onto the space of tensions $\operatorname{im}N$, and let $C=I-P$, the projection onto the space of circulations $\ker N^{\top}$ of the introduction.

\begin{lem}\label{lem:energy}
For every orientation $s$, $E(s)=\|Ps\|^2=m-\|Cs\|^2$.
\end{lem}

\begin{proof}
With $f=\Delta^{+}N^{\top}s$ we have $\Delta f=N^{\top}s$ and $E(s)=\langle f,N^{\top}s\rangle=s^{\top}N\Delta^{+}N^{\top}s=\|Ps\|^2$. Since $\|s\|^2=m$ and $P+C=I$ with orthogonal ranges, $\|Ps\|^2=m-\|Cs\|^2$.
\end{proof}

\begin{lem}\label{lem:strong}
Let $G$ be bridgeless. The least energy over all orientations of $G$ is attained by a strongly connected orientation; hence $\Emin(G)=m-\max_{s\in\{\pm1\}^E}\|Cs\|^2$. Moreover, every maximizer $s$ of $\|Cs\|^2$ is a strongly connected orientation and satisfies $s_e(Cs)_e\ge C_{ee}>0$ for every edge $e$.
\end{lem}

\begin{proof}[Proof of Lemma~\ref{lem:strong}]
Let $s$ maximize $s^{\top}Cs$ over $\{\pm1\}^E$ and put $x=Cs$. Changing the sign of $s_e$ changes $s^{\top}Cs$ by $-4s_ex_e+4C_{ee}$, so $s_ex_e\ge C_{ee}$. Each edge $e$ lies on a cycle; if $\gamma$ is the signed indicator vector of that cycle, then $\gamma/\sqrt{|\gamma|}$ is a unit circulation with coordinate $\pm1/\sqrt{|\gamma|}$ at $e$, where $|\gamma|\le m$ is the length of the cycle. Hence $C_{ee}=\|C\mathbf 1_e\|^2\ge1/|\gamma|>0$. Thus $x$ is a circulation with $s_ex_e\ge C_{ee}>0$ for every edge, that is, a circulation that is positive on every arc of the orientation $s$. If that orientation were not strongly connected, some strongly connected component $X$ would have arcs leaving it and none entering it, and the net flow of $x$ out of $X$ would be positive, which is impossible for a circulation. By Lemma~\ref{lem:energy}, $s$ minimizes the energy over all orientations.
\end{proof}

The support of a nonzero circulation has no vertex of degree one, so it contains a cycle; the signed indicator vector of a cycle has minimal support among nonzero circulations, because the support of a nonzero circulation is never a forest, so it is elementary in $\Lambda(G)$. So $\Lambda(G)$ is regular, and the theorem of Dutour Sikiri\'c, Sch\"urmann and Vallentin \cite[Theorem~2]{DSV09} shows that its Voronoi cell is $C([-1/2,1/2]^E)$, with vertices among the points $Cs/2$. For a strongly connected orientation $D$ with sign vector $s_D$, the point $Cs_D/2$ lies in the space of circulations and satisfies $2\langle Cs_D/2,\gamma\rangle=\langle s_D,\gamma\rangle=|\gamma|$ for the signed indicator vector $\gamma$ of every directed cycle of $D$, where $|\gamma|$ is the length of the cycle, which are the equations defining $v^D$ in \cite[Lemma~8]{Amini}; so $v^D=Cs_D/2$. The identity of Lemma~\ref{lem:strong} and Corollary~\ref{cor:radius} below is that of \cite[Theorem~2]{DSV09} together with \cite[Lemma~8, Proposition~12]{Amini}.

\begin{cor}\label{cor:radius}
The squared covering radius of $\Lambda(G)$ is the largest of the numbers $\|Cs\|^2/4$ with $s\in\{\pm1\}^E$. If $G$ is bridgeless with $m$ edges, it equals $(m-\Emin(G))/4$.
\end{cor}

\begin{proof}
The cell $C([-1/2,1/2]^E)$ contains every point $Cs/2$ and its vertices are among them, so the squared covering radius, the largest squared norm of a point of the cell, is $\max_s\|Cs\|^2/4$. For bridgeless $G$ this is $(m-\Emin(G))/4$ by Lemma~\ref{lem:strong}.
\end{proof}

In the proof of \cite[Lemma~13]{Amini} the coordinate $v^D_e$ on an arc $e=(y,z)$ of $D$ is written $f_D(z)-f_D(y)+1$; with the value $(f_D(z)-f_D(y)+1)/2$ the same computation gives $\|v^D\|^2=(m-E(D))/4$, so \cite[Definition--Theorem~14]{Amini} holds with $1/4$ in place of $1/2$. For a cycle of length $n$ the value $(m-\Emin(G))/4$ is $n/4$, the squared distance from $\tfrac12(1,\dots,1)$ to the lattice $\Z(1,\dots,1)$.

\begin{lem}\label{lem:lower}
Let $G$ be bridgeless and let $v$ be a vertex of odd degree. Then $\Emin(G)\ge1/\deg(v)$.
\end{lem}

\begin{proof}
Since $N^{\top}C=0$, $N^{\top}Ps=N^{\top}s$, and $(N^{\top}s)_v=d^{\mathrm{out}}_v-d^{\mathrm{in}}_v$ is a nonzero integer when $\deg(v)$ is odd. Its absolute value is at most $\sum_{e\ni v}|(Ps)_e|$, so by the Cauchy--Schwarz inequality $\|Ps\|^2\ge1/\deg(v)$.
\end{proof}

If every degree is even, orienting the edges along a closed walk that uses each edge once gives a strongly connected orientation with $d^{\mathrm{out}}=d^{\mathrm{in}}$, and $\Emin(G)=0$.

\section{Hardness}\label{sec:hard}

The reduction has three steps. A not-all-equal formula is turned into a sum of fractions that can be signed to an integer exactly when the formula is satisfiable (Lemma~\ref{lem:signed}). Each fraction becomes a weighted triangle in the ring $G_0$, and the graph $H$ is a subdivided plane dual of $G_0$. Finally the least energy of $H$ is computed in closed form (Proposition~\ref{prop:ring}).

\subsection{A signed sum}
We reduce from monotone not-all-equal 3-satisfiability with exactly four occurrences per variable, \textsc{Monotone NAE-3-Sat-E4}: given clauses, each consisting of three distinct variables, with every variable in exactly four clauses and no negations, is there an assignment of the values $\pm1$ such that no clause has all three values equal? Darmann and D\"ocker proved that this problem is NP-complete \cite[Theorem~1]{DD20}.

Let the clauses be $c_1,\dots,c_q$. For each clause $c_j$ add a new variable $w_j$ occurring only in $c_j$. Three signs are not all equal exactly when a fourth sign can be chosen so that the four sum to zero, so the instance is a yes-instance if and only if there are signs for all variables such that every extended clause $c_j\cup\{w_j\}$ has sign sum $0$. Let $p_1<\dots<p_q$ be the first $q$ primes larger than $10$; they are at most the $(q+4)$-th prime, which is below $n(\log n+\log\log n)<n^2$ for $n=q+4\ge8$, with natural logarithms \cite{RS62}; here $q\ge4$, since every variable lies in four clauses. So each has at most $2\log_2(q+4)+1$ bits, and they can be listed by trial division in time polynomial in $q$. For every variable $i$ (original or new) put
\[
a_i=\sum_{j:\,i\in c_j\cup\{w_j\}}\frac1{p_j}=\frac{A_i}{B_i}\quad\text{in lowest terms}.
\]
Each $i$ lies in at most four extended clauses, so $0<a_i\le 4/11<1/2$, the denominator $B_i$ is odd and $B_i\le p_q^4$. With $Q=p_1\cdots p_q$, for these $a_i$ the quantity $\rho$ of the introduction is
\[
\rho=\min_{\sigma}\dist\Bigl(\sum_i\sigma_ia_i,\ \Z\Bigr).
\]

\begin{lem}\label{lem:signed}
$\rho=0$ if the instance is a yes-instance, and $\rho\ge1/Q$ otherwise.
\end{lem}

\begin{proof}
Let $\tau_j(\sigma)$ be the sign sum of the extended clause $c_j\cup\{w_j\}$, so $|\tau_j|\le4$ and $\sum_i\sigma_ia_i=\sum_j\tau_j/p_j$. If every $\tau_j$ vanishes the sum is $0$. Conversely $Q\sum_i\sigma_ia_i=\sum_j\tau_jQ/p_j$ is an integer. If $\sum_i\sigma_ia_i$ is an integer $\nu$, reducing $Q\nu=\sum_l\tau_lQ/p_l$ modulo $p_j$ gives $\tau_jQ/p_j\equiv0$, so $p_j$ divides $\tau_j$, and $|\tau_j|\le4<p_j$ forces $\tau_j=0$ for all $j$. Hence if the instance is a no-instance, $\sum_i\sigma_ia_i$ is never an integer, and as an integer multiple of $1/Q$ it lies at distance at least $1/Q$ from $\Z$.
\end{proof}

\subsection{The graph}
Let $k$ be the number of variables, original and new, indexed cyclically by $i\in\Z/k\Z$. Build the ring $G_0$, a weighted plane graph: take vertices $u_1,\dots,u_k$ on a circle and, for each $i$, a triangle on $u_i,u_{i+1}$ and a new vertex $u'_i$, with the edge $u_iu_{i+1}$ (the \emph{inner} edge) of weight $2A_i$ and the edges $u_iu'_i$, $u'_iu_{i+1}$ (the \emph{outer} edges) of weight $2(B_i-2A_i)$ each. Draw the inner edges around a central face and the outer edges around the outer face. The faces of $G_0$ are the $k$ triangles, the central face and the outer face.

Let $H$ be obtained from the plane dual of $G_0$ by replacing each dual edge, dual to a primal edge of weight $w$, by a path of $w$ edges. So $H$ has a vertex $o_{\mathrm{in}}$ for the central face, a vertex $o_{\mathrm{out}}$ for the outer face, a vertex $h_i$ for each triangle, a path of length $2A_i$ from $h_i$ to $o_{\mathrm{in}}$ and two paths of length $2(B_i-2A_i)$ from $h_i$ to $o_{\mathrm{out}}$. All weights are positive even integers bounded by $2p_q^4$, so $H$ has polynomial size; here the bounded number of occurrences matters.

The graph $H$ is planar, as a subdivision of a plane dual. It is simple because every path has length at least two. It is bipartite because every cycle of $H$ is a union of whole paths, each of even length. It is bridgeless: the two paths from $h_i$ to $o_{\mathrm{out}}$ form a cycle, and the path from $h_i$ to $o_{\mathrm{in}}$ lies on the cycle through $h_i,o_{\mathrm{in}},h_{i+1},o_{\mathrm{out}}$. Deleting $o_{\mathrm{out}}$ leaves a subdivided star at $o_{\mathrm{in}}$ with pendant paths attached, which is a tree; adding $o_{\mathrm{out}}$ to every bag of a width-one tree decomposition of this tree gives treewidth at most two. The number of edges of $H$ is $W=\sum_i\bigl(2A_i+4(B_i-2A_i)\bigr)$.

\subsection{The least energy of \texorpdfstring{$H$}{H}}
A circulation of $H$ is constant along each subdivided path, so circulations of $H$ correspond to circulations of the plane dual of $G_0$, and these correspond to tensions of $G_0$, that is, to vectors $(y_a-y_b)_{ab\in E(G_0)}$ for potentials $y\colon V(G_0)\to\R$ \cite[Chapter~14]{GR01}. A circulation corresponding to the tension $\delta$ has squared norm $\sum_ew_e\delta_e^2$ and $\ell^1$ norm $\sum_ew_e|\delta_e|$, where $w_e$ is the weight of $e$. By Lemma~\ref{lem:strong}, and since $\max_s\|Cs\|^2=\max_s\max_{\varphi}\bigl(2\langle s,\varphi\rangle-\|\varphi\|^2\bigr)=\max_\varphi\bigl(2\|\varphi\|_1-\|\varphi\|^2\bigr)$ with $\varphi$ ranging over circulations,
\begin{equation}\label{eq:dual}
\Emin(H)=W-4R^2,\qquad R^2=\max_{\theta\colon V(G_0)\to\R}\ \sum_{ab\in E(G_0)}w_{ab}\,g(\theta_a-\theta_b),
\end{equation}
where $g(\delta)=|\delta|-\delta^2$, after the substitution $y=2\theta$. By Corollary~\ref{cor:radius}, $R^2$ is the squared covering radius of $\Lambda(H)$.

Let $\hat g(\delta)=\{\delta\}(1-\{\delta\})$, where $\{\delta\}$ is the fractional part. Then $\hat g$ has period $1$, $\hat g=g$ on $(-1,1)$, and $g<0\le\hat g$ outside $[-1,1]$. Hence $g\le\hat g$, and for $\theta$ with values in $[0,1)$ every difference lies in $(-1,1)$. So $R^2$ is also the maximum of $\sum w_{ab}\hat g(x_a-x_b)$, a function of $\theta$ modulo $1$.

\begin{lem}\label{lem:triangle}
Let $0<A<B/2$ be integers, $a=A/B$ and $F=B-2A+A^2/B$. For $t\in\R$,
\[
\max_{\zeta\in\R}\Bigl(2A\,\hat g(t)+2(B-2A)\bigl(\hat g(\zeta)+\hat g(t-\zeta)\bigr)\Bigr)=F-B\,\dist\bigl(t,\{a,-a\}+\Z\bigr)^2 .
\]
\end{lem}

\begin{proof}
Both sides have period $1$ and are even in $t$, so let $t\in[0,1/2]$. For $\alpha=\{\zeta\}$ and $\beta=\{t-\zeta\}$ we have $\alpha+\beta\in\{t,t+1\}$ and $\hat g(\zeta)+\hat g(t-\zeta)=(\alpha+\beta)-(\alpha^2+\beta^2)$, which for a fixed sum $\kappa=\alpha+\beta$ is at most $\kappa-\kappa^2/2$, with equality at $\alpha=\beta=\kappa/2$. Taking $\kappa=t$ or $\kappa=t+1$, the largest value is $\max(t,1/2)-t^2/2=1/2-t^2/2$ on $[0,1/2]$. The left side is therefore
\[
2A(t-t^2)+(B-2A)(1-t^2)=(B-2A)+2At-Bt^2=F-B(t-a)^2 .
\]
For $t\in[0,1/2]$ the nearest point of $\{a,-a\}+\Z$ is $a$, since $a<1/2$.
\end{proof}

Let $F_i=B_i-2A_i+A_i^2/B_i$, so that $U=\sum_iF_i$, with $U$ and $S$ as in the introduction.

\begin{prop}\label{prop:ring}
$R^2=U-\rho^2/S$, hence $\Emin(H)=W-4U+4\rho^2/S$.
\end{prop}

\begin{proof}
Maximize over the potentials of the $u'_i$ first. By Lemma~\ref{lem:triangle}, with $t_i=\theta_{u_{i+1}}-\theta_{u_i}$,
\[
R^2=U-\min\Bigl\{\sum_iB_i\dist(t_i,\{\pm a_i\}+\Z)^2:\ \sum_it_i\in\Z\Bigr\},
\]
because the differences $t_i$ around the circle $u_1,\dots,u_k$ can be any real numbers with integer sum (working modulo $1$). Write $t_i=\sigma_ia_i+n_i+\eta_i$ with $\sigma_i=\pm1$, $n_i\in\Z$ and $|\eta_i|$ the distance. Then $\sum_i\eta_i\equiv-\sum_i\sigma_ia_i\pmod1$, so $|\sum_i\eta_i|\ge \rho$, and by the Cauchy--Schwarz inequality $\bigl(\sum_i\eta_i\bigr)^2\le S\sum_iB_i\eta_i^2$. Thus the minimum is at least $\rho^2/S$. Choosing $\sigma$ that attains $\rho$, an integer $\nu$ at distance $\rho$ from $\sum\sigma_ia_i$, and $\eta_i=(\nu-\sum_j\sigma_ja_j)/(B_iS)$ gives $\sum_iB_i\eta_i^2=\rho^2/S$ with $\sum_it_i\in\Z$, and the distance from $t_i$ to $\{\pm a_i\}+\Z$ is at most $|\eta_i|$. So the minimum is $\rho^2/S$. The formula for $\Emin(H)$ is \eqref{eq:dual}.
\end{proof}

\begin{proof}[Proof of Theorem~\ref{thm:hard}]
The problem is in NP: a strongly connected orientation can be checked by two searches, and its energy is computed exactly by solving a rational linear system whose solution has polynomially many bits. For hardness, map an instance of \textsc{Monotone NAE-3-Sat-E4} to $H$ and $r_0=W-4U$, a rational number with polynomially many bits. By Proposition~\ref{prop:ring} and Lemma~\ref{lem:signed}, $\Emin(H)=r_0$ for a yes-instance and $\Emin(H)\ge r_0+4/(Q^2S)>r_0$ for a no-instance.
\end{proof}

\begin{proof}[Proof of Corollary~\ref{cor:cov}]
By Corollary~\ref{cor:radius} the squared covering radius of $\Lambda(G)$ is the largest of the numbers $\|Cs\|^2/4$. A sign vector $s$ with $\|Cs\|^2/4>r$ certifies a no-instance and is checked by one rational linear solve, so the problem is in coNP. For the instance $(H,r_0)$ of the proof of Theorem~\ref{thm:hard}, put $r=(W-r_0-2/(Q^2S))/4$. The squared covering radius of $\Lambda(H)$ is at most $r$ exactly when $\Emin(H)\ge r_0+2/(Q^2S)$, that is, for a no-instance of the formula. This reduces the complement of an NP-complete problem to ours. For a plane multigraph $G$ the circulation space of $G$ is the tension space of its plane dual $G^*$ under the natural bijection of edges \cite[Chapter~14]{GR01}, so $\Lambda(G)$ is isometric to the lattice of integer cuts of $G^*$. Lattices of integer cuts are regular \cite[Lemma~1(i)]{DSV09}, so by \cite[Theorem~2]{DSV09} the squared covering radius of the lattice of integer cuts of $G$ is the largest of the numbers $\|Ps\|^2/4$, and the same certificate shows that the problem for them is in coNP. Taking $G=H$, whose plane dual $H^*$ is a planar multigraph, the lattice of integer cuts of $H^*$ is isometric to $\Lambda(H)$, so the problem is coNP-hard for lattices of integer cuts of planar multigraphs. This is the hardness anticipated after \cite[Proposition~3]{DSV09}.
\end{proof}

\begin{ex}\label{ex:small}
For $(A_i,B_i)=(1,3),(1,3),(1,3)$ the graph $H$ has $W=18$ edges, $U=4$ and $\rho=0$, and $\Emin(H)=2$. For $(1,3),(1,3),(1,5)$ it has $W=26$, $U=88/15$, $S=13/15$ and $\rho=2/15$, and $\Emin(H)=34/13$. A strongly connected orientation of $H$ is constant on each subdivided path, since an inner vertex of a path would otherwise be a source or a sink.
\end{ex}

\section{Approximation at bounded treewidth}\label{sec:fptas}

\begin{proof}[Proof of Theorem~\ref{thm:fptas}]
If every degree is even, output an Eulerian orientation with $\ell=0$. Otherwise let $M=\Emin(G)$, so $M\ge1/m$ by Lemma~\ref{lem:lower}. Let $s$ and $x=Cs$ be as in the proof of Lemma~\ref{lem:strong}, so $s_ex_e\ge C_{ee}\ge1/m$ and $\|s-x\|^2=\|Ps\|^2=M$. Fix an integer $K>m$ with $K^2\varepsilon\ge m^2$, for instance $K=1+\lceil m/\sqrt\varepsilon\rceil$.

The vector $Kx$ is a circulation, and since the incidence matrix is totally unimodular, the Hoffman--Kruskal theorem gives an integral circulation $z$ with $\lfloor Kx_e\rfloor\le z_e\le\lceil Kx_e\rceil$ for every $e$ \cite[Chapter~19]{Schrijver}. As $|Kx_e|\ge K/m>1$, every $z_e$ is nonzero with the sign $s_e$, and $|z_e|\le T=\lceil K\sqrt m\rceil$ because $\|x\|\le\|s\|=\sqrt m$. Consider
\[
J=\min\Bigl\{\sum_e\bigl(|z'_e|/K-1\bigr)^2:\ z'\in\Z^E\ \text{a circulation},\ 1\le|z'_e|\le T\Bigr\}.
\]
For a feasible $z'$ let $s'=\operatorname{sign}(z')$. Then $\sum_e(|z'_e|/K-1)^2=\|s'-z'/K\|^2\ge\|Ps'\|^2=E(s')$, because $z'/K$ is a circulation; and $s'$ is strongly connected, because $z'$ is a circulation positive on every arc of $s'$ (as in the proof of Lemma~\ref{lem:strong}). For the rounded $z$, the vector $s-x$ is a tension and $x-z/K$ a circulation, so
\[
\|s-z/K\|^2=\|s-x\|^2+\|x-z/K\|^2\le M+m/K^2\le(1+\varepsilon)M,
\]
using $m/K^2\le\varepsilon/m\le\varepsilon M$. Hence an optimal $z'$ for $J$ gives a strongly connected orientation $D$ with $M\le E(D)\le J\le(1+\varepsilon)M$. Put $\ell=\max(1/m,\,J-m/K^2)$. Then $J-m/K^2\le\|s-z/K\|^2-m/K^2\le M$ and $J-m/K^2\ge M-\varepsilon M$, so $(1-\varepsilon)M\le\ell\le M$. The energy $E(D)$ is computed exactly by a rational linear solve.

It remains to compute $J$ and an optimal $z'$. Compute a tree decomposition of width at most $\omega$ in linear time \cite{Bodlaender96}. Assign every edge to one bag containing both ends, and process the decomposition from the leaves, keeping for each bag a table indexed by the vector of net flows out of its vertices from the edges assigned so far, each entry the least partial cost $\sum(|z'_e|-K)^2$ together with a witness. An edge is introduced by trying the $2T$ values of $z'_e$; a vertex is forgotten only with net flow $0$; the tables of the children of a bag are combined one at a time by adding the vectors of net flows and the costs. Every net flow lies in $[-mT,mT]$, so each table has at most $(2mT+1)^{\omega+1}$ entries and each step takes time polynomial in $m$ and $(2mT+1)^{\omega+1}$. At the root the entry with all net flows $0$ is $K^2J$. Since $T\le 2m^{3/2}/\sqrt\varepsilon+3\sqrt m$, the total time is polynomial in $m$ and $1/\varepsilon$ for fixed $\omega$.
\end{proof}

In Theorem~\ref{thm:hard} the gap is $4/(Q^2S)$, while $r_0=W-4U=\sum_i2A_i(1-2A_i/B_i)$ lies between $6k/11$ and $W$; the guarantee of Theorem~\ref{thm:fptas} separates the two cases only when $\varepsilon<4/(Q^2Sr_0)$, which is exponentially small in the size of the formula.

\section{Open questions}\label{sec:open}

\begin{question}
Is $\Emin$ NP-hard to approximate within some constant factor on general bridgeless graphs, or does it admit a polynomial-time approximation scheme? Is it hard to compute exactly when the threshold $r$ is required to have a denominator bounded by a polynomial in the number of vertices?
\end{question}

\begin{question}
Theorem~\ref{thm:fptas} needs bounded treewidth. Is there a polynomial-time approximation scheme for $\Emin$ on planar graphs?
\end{question}

\begin{question}
Amini's description of the Voronoi cell was extended to the lattice of integer flows of a regular matroid by Dancso and Lim~\cite{DL22}. For which classes of regular matroids is the covering radius of the lattice of integer flows computable exactly in polynomial time? Matroids of bounded branch width are a natural first case.
\end{question}

\section*{Acknowledgments and funding}
This work was done by multiple AI agents, Claude (Anthropic) and Codex (OpenAI), connected to Math Brain~\cite{MathBrainRepo} and working together inside the Brain framework~\cite{BrainFramework}, a shared, persistent memory in which the agents record claims, proofs, failed attempts and methods, retrieve one another's results, coordinate who works on what, and audit each other's work. The method is described in~\cite{MathBrainPaper}.

The research presented in this paper was supported by the European Research Council (ERC) under the European Union's Horizon Europe research and innovation program (grant agreement No.~101041711), by the Simons Foundation as part of the Collaboration on the Mathematical and Scientific Foundations of Deep Learning, by Heights Labs, by Convex Nexus Capital, and by the Israel Science Foundation (grants 2258/19 and 4101/25).

\section*{Declaration of competing interest}
The author holds equity in Heights Labs and in Convex Nexus Capital and has a management role at Convex Nexus Capital.

\section*{Data availability}
The program \texttt{check\_triangle\_ring.py}, which uses only Python's standard library, computes the values of Example~\ref{ex:small} from the definition, in exact rational arithmetic, over the orientations that are constant on each subdivided path; for the first graph of Example~\ref{ex:small} it also runs over all $2^{18}$ orientations. It and its output are in the repository
\begin{center}\url{https://github.com/yspennstate/minimum-energy-strong-orientations/tree/v1.3}.\end{center}

\section*{Declaration of generative AI and AI-assisted technologies}
The AI agents Claude (Anthropic) and Codex (OpenAI), working through Math Brain~\cite{MathBrainRepo,MathBrainPaper} on the Brain framework~\cite{BrainFramework}, produced the statements, proofs and computations of this paper under the author's direction. The question is Amini's~\cite{Amini}. The description of the Voronoi cell through strongly connected orientations is his, and the description of Voronoi cells of regular lattices as projected cubes is due to Dutour Sikiri\'c, Sch\"urmann and Vallentin~\cite{DSV09}. The reduction starts from the theorem of Darmann and D\"ocker~\cite{DD20}. The lattice of integer flows was studied earlier by Bacher, de la Harpe and Nagnibeda~\cite{BHN97}. The description of its Voronoi cell was conjectured by Caporaso and Viviani~\cite{CV10} and extended to regular matroids by Dancso and Lim~\cite{DL22}. The hardness of the covering radius problem for general lattices is due to Vallentin~\cite{Vallentin26}. The hardness results for acyclic orientations recalled in the introduction are due to Borradaile, Iglesias, Migler, Ochoa, Wilfong and Zhang~\cite{BIMOWZ17} and to Borsik and Madarasi~\cite{BM25}, and the polynomial-time algorithm for the closest vector problem recalled there is due to McCormick, Peis, Scheidweiler and Vallentin~\cite{MPSV21}. The values of Example~\ref{ex:small} were computed with the program named in the data availability statement. The author directed the work and is responsible for the content of the paper and for decisions concerning its dissemination.

\begin{thebibliography}{99}
\bibitem{Amini} O.~Amini, Lattice of integer flows and poset of strongly connected orientations, preprint, 2010, revised 2020, arXiv:1007.2456.
\bibitem{BHN97} R.~Bacher, P.~de la Harpe and T.~Nagnibeda, The lattice of integral flows and the lattice of integral cuts on a finite graph, \emph{Bull. Soc. Math. France} \textbf{125} (1997), 167--198, \doi{10.24033/bsmf.2303}.
\bibitem{Bodlaender96} H.~L. Bodlaender, A linear-time algorithm for finding tree-decompositions of small treewidth, \emph{SIAM J. Comput.} \textbf{25} (1996), 1305--1317, \doi{10.1137/S0097539793251219}.
\bibitem{BIMOWZ17} G.~Borradaile, J.~Iglesias, T.~Migler, A.~Ochoa, G.~Wilfong and L.~Zhang, Egalitarian graph orientations, \emph{J. Graph Algorithms Appl.} \textbf{21} (2017), 687--708, \doi{10.7155/jgaa.00435}.
\bibitem{BM25} N.~A. Borsik and P.~Madarasi, Separable convex optimization over indegree polytopes, preprint, 2025, arXiv:2509.06182.
\bibitem{CV10} L.~Caporaso and F.~Viviani, Torelli theorem for graphs and tropical curves, \emph{Duke Math. J.} \textbf{153} (2010), 129--171, \doi{10.1215/00127094-2010-022}.
\bibitem{DL22} Z.~Dancso and J.~Lim, Lattice of integer flows and the poset of strongly connected orientations for regular matroids, \emph{Discrete Math.} \textbf{346} (2023), 113589, \doi{10.1016/j.disc.2023.113589}; arXiv:2202.13265.
\bibitem{DD20} A.~Darmann and J.~D\"ocker, On a simple hard variant of Not-All-Equal 3-Sat, \emph{Theoret. Comput. Sci.} \textbf{815} (2020), 147--152, \doi{10.1016/j.tcs.2020.02.010}.
\bibitem{DSV09} M.~Dutour Sikiri\'c, A.~Sch\"urmann and F.~Vallentin, Complexity and algorithms for computing Voronoi cells of lattices, \emph{Math. Comp.} \textbf{78} (2009), 1713--1731, \doi{10.1090/S0025-5718-09-02224-8}.
\bibitem{GR01} C.~Godsil and G.~Royle, \emph{Algebraic Graph Theory}, Graduate Texts in Mathematics 207, Springer, New York, 2001.
\bibitem{MPSV21} S.~T. McCormick, B.~Peis, R.~Scheidweiler and F.~Vallentin, A polynomial time algorithm for solving the closest vector problem in zonotopal lattices, \emph{SIAM J. Discrete Math.} \textbf{35} (2021), 2345--2356, \doi{10.1137/20M1382258}.
\bibitem{Robbins39} H.~E. Robbins, A theorem on graphs, with an application to a problem of traffic control, \emph{Amer. Math. Monthly} \textbf{46} (1939), 281--283, \doi{10.2307/2303897}.
\bibitem{RS62} J.~B. Rosser and L.~Schoenfeld, Approximate formulas for some functions of prime numbers, \emph{Illinois J. Math.} \textbf{6} (1962), 64--94, \doi{10.1215/ijm/1255631807}.
\bibitem{Schrijver} A.~Schrijver, \emph{Theory of Linear and Integer Programming}, Wiley, Chichester, 1986.
\bibitem{Vallentin26} F.~Vallentin, The complexity of computing the covering radius of a Euclidean lattice, preprint, 2026, arXiv:2609.35027.
\bibitem{MathBrainRepo} Y.~Shmalo, \emph{Math Brain: public snapshot, code and tests}, 2026, \url{https://github.com/yspennstate/math-brain-public}.
\bibitem{MathBrainPaper} Y.~Shmalo, \emph{Math Brain: shared, auditable memory for multi-agent mathematical research}, October 2026, \url{https://github.com/yspennstate/math-brain-public/tree/main/papers/math-brain-2026-10-24}, \doi{10.5281/zenodo.24150837}.
\bibitem{BrainFramework} Y.~Shmalo, \emph{Brain Framework}, public software repository, 2026, \url{https://github.com/yspennstate/brain-framework-public}.
\end{thebibliography}
\end{document}
